\uuid{5zea}
\titre{Fraction}
\auteur{bodin ; hanani ; mzari}
\organisation{exo7}
\licence{CC BY-NC-SA 4.0}
\datecreate{2020-01-20}
\niveau{L1}
\module{Analyse}
\chapitre{Propriétés de R}
\sousChapitre{Les rationnels}
\difficulte{1}
\theme{Rationnels, fraction}
\exo7qcmid{lille-1.reels.002}

\contenu{
\begin{question}
\qclassification{120.01}
\qtags{motcle=fraction}
Quelles sont les assertions vraies ?
\begin{answers}
    \bad{$\frac{1}{7} = 0,142142142\ldots$}

    \good{Le nombre dont l'écriture décimale est $0,090909\ldots$ est un nombre rationnel.}

    \good{$9,99999\ldots = 10$}

    \bad{$\frac{1}{5} = 0,202020\ldots$}

\end{answers}
\begin{explanations}
On trouve l'écriture décimale d'un rationnel en calculant la division !
Si on a une écriture décimale finie ou périodique alors c'est un nombre rationnel.
\end{explanations}
\end{question}
}
